\chapter{Дифференцируемый рендеринг разреженных SDF}
\label{ch:chap2}

Базовый метод дифференцируемого физически-обоснованного рендеринга \cite{Relax2024} предлагает простую оценку граничного интеграла для функций расстояния со знаком -- более гибкого представления с точки зрения реконструкции поверхности, чем меш. Тем не менее, у него имеется ряд недостатков; во-первых, быстрый рост размера модели из-за использования регулярной сетки: в процессе оптимизации разрешение сетки увеличивается в 8 раз после каждой эпохи, что сильно ограничивает число эпох, которые можно выполнить при небольшом объёме памяти. Во-вторых, медленная отрисовка, основанная на классическом методе поиска пересечения sphere tracing. В-третьих, даже при локальных изменениях дистанций около поверхности, существующие методы применяют затратный алгоритм редистансинга. Он применяется после каждого шага оптимизации для пересчёта значений, так как сетка перестаёт задавать функцию расстояния, причём пересчитываются все дистанции сетки.

Предлагается метод дифференцируемого рендеринга разреженных SDF представлений. Взяв метод релаксированной границы в качестве базового, он предлагает переработанные этапы процесса оптимизации, а также включает новые алгоритмы, которые направлены на поддержку разреженных SDF представлений (см. \hyperref[fig:comp]{Рисунок \ref*{fig:comp}}). Ключевой вклад метода заключается в следующем:

\begin{figure}[ht]
    \centering
    \includegraphics[width=\textwidth]{header.eps}
    \caption{Существующие методы (слева) используют регулярную SDF сетку для представления поверхности, sphere tracing для пересечения лучей со сценой и равномерную выборку изображения при проецировании лучей. Предложенный метод (справа) использует разреженное множество бриков (маленьких SDF сеток), где луч пересекает каждый брик на своем пути, пока не достигнет поверхности. Он также использует адаптивную выборку для граничного интеграла, пуская больше лучей в области релаксированной границы. Красный и синий цвета отображают положительные и отрицательные значения SDF соответственно}
    \label{fig:comp}
\end{figure}

\begin{enumerate}
    \item Адаптация дифференцируемого рендеринга SDF под разреженные представления, превосходящая базовый метод по скорости и уменьшающая на порядок рост размера оптимизируемой модели.
    \item Две версии алгоритма редистансинга разреженных представлений, одна из которых адаптирована непосредственно под нужды дифференцируемого рендеринга.
\end{enumerate}

Реконструкция происходит в несколько эпох. Каждая эпоха состоит из нескольких сотен итераций, каждая итерация состоит из следующих шагов: сперва проводится отрисовка группы изображений с референсных ракурсов, в процессе отрисовки аккумулируются градиенты; после этого к дистанциям применяется регуляризация; расчитывается функция потерь между полученными изображениями и соответствующими им референсами, следом происходит обратное распространение ошибки, в конце итерации применяется локальный редистансинг. Между эпохами SDF сетка разреживается, повышается её разрешение, после чего к ней применяется глобальный редистансинг.

Метод использует в реконструкции представление, называемое разреженным множеством бриков (sparse brick set, SBS) \cite{HanssonSoderlund}. Это гибрид между иерархическим и регулярным представлением для функции расстояния со знаком; он представляет собой BVH-дерево, в листьях которого содержатся брики -- маленькие регулярные SDF сетки. Благодаря этому можно игнорировать пустое пространство, отбрасывая те листья, в которых отсутствует поверхность объекта. Сперва луч обходит BVH-дерево; если найдено пересечение с листом, луч пересекается с содержащейся в нём SDF сеткой. Луч пересекает каждый лист на своём пути, пока не попадёт в поверхность или на его пути не останется ни одного листа. Помимо самого представления, в нашем методе используются предложенные для SBS методы поиска пересечения луча: метод Ньютона и аналитический метод, адаптированные под поиск точек релаксированной границы для корректного расчёта градиентов. Оба метода пересекают луч с вокселем и основаны на том, что внутри вокселя, SDF вдоль луча представима в виде кубического полинома, отличаются эти методы только способом нахождения корня. В то же время, основной критерий точки релаксированной границы: эта точка является локальным минимумом SDF вдоль луча, поэтому адаптация поиска точек релаксированной границы внутри вокселя сводится к нахождению нужного корня квадратного уравнения. Однако, при использовании трилинейной интерполяции распространённой является ситуация, когда локальный минимум достигается на границе вокселей -- этот случай также обрабатывается.

Алгоритм реализован с использованием C++ and Vulkan, без систем автоматического дифференцирования -- все градиенты получены аналитически. Для поддержки произвольного числа семплов на пиксель в шейдере, используемом для отрисовки оптимизируемой модели и оценки граничного интеграла, реализовано аккумулирование градиентов каждые 16 семплов. Используется цвет, который удалось аккумулировать на данный момент, а не финальный цвет пикселя. Это нужно, так как для каждого семпла требуется хранить информацию о том, в какой воксель он попал, а для точной оценки градиента требуется цвет пикселя, полученный после семплирования; при этом число семплов задаётся пользователем и заранее неизвестно. Мы исходим из предположения, что 16 семплов достаточно для получения цвета, близкого к финальному, и храним информацию о вокселях в стеке фиксированного размера.

% Редистансинг

\begin{figure}[ht]
    \centering
    \includegraphics[width=0.9\textwidth]{redistancing.eps}
    \caption{Сравнение глобального алгоритма редистансинга на плотной сетке (слева) с локальным редистансингом для SBS (справа); размеры бриков -- 3x3 дистанции. Точки представляют дистанции. Красные и синие точки -- дистанции вне и внутри объекта соответственно. Пунктирная линия представляет поверхность, заданную дистанциями белого цвета. Стрелки показывают порядок обновления значений. Локальный редистансинг игнорирует брики без поверхности (показаны серым)}
    \label{fig:redist}
\end{figure}

После шага градиентного спуска обновлённые дистанции, строго говоря, больше не задают SDF: они не соответствуют настоящим расстояниям до поверхности, и сетка больше не удовлетворяет уравнению Эйконала. Для восстановления этого свойства в базовом методе применяется реализация алгоритма редистансинга, называемая методом fast sweeping \cite{RedistFSM}. Адаптация этого шага под разреженные представления нетривиальна: при пересчёте дистанций в одном брике модели нужно дополнительно найти и учитывать соседние брики, не связанные с ним дистанциями -- случай, когда промежуточные брики удалены, так как не содержат поверхность, но сам сосед влияет на корректное значение дистанций.

Полноценной адаптацией редистансинга под SBS является глобальный редистансинг. Для корректного расчёта дистанций добавлен шаг экстраполяции: для каждой грани брика, к которой не примыкает другой брик (нет непосредственного соседа), ищутся соседи, и дистанции грани экстраполируются из ближней грани этих соседей. Выбираются минимальные по модулю значения дистанций среди текущих и всех экстраполированных. Сам редистансинг распараллелен по брикам, поскольку существующие схемы параллелизации для метода fast sweeping основаны на самих значениях, что алгоритмически сложно реализовать для SBS.

Одним полезным свойством аналитического метода и метода Ньютона для пересечения луча с вокселем является меньшая зависимость от корректных дистанций, чем у sphere tracing. Благодаря этому можно игнорировать воксель, если все его дистанции больше, чем relaxed epsilon: в этом случае мы не найдём в нём ни пересечение, ни точку релаксированной границы. В связи с этим нас меньше интересует значение дистанций в бриках, которые не хранят поверхность, в частности, к ним необязательно применять редистансинг. Локальный редистансинг является оптимизированной под дифференцируемый рендеринг версией шага редистансинга, применяемой после каждой итерации метода (см. \hyperref[fig:redist]{Рисунок \ref*{fig:redist}}). Первая оптимизация заключается в том, что пересчитываются дистанции только в бриках, содержащих поверхность. Вторая оптимизация: предполагается, что дистанции в бриках с поверхностью зависят только от поверхности в этом брике -- не происходит обращения к дистанциям в других бриках, что упрощает логику и увеличивает локальность без ухудшения качества реконструкции.

Для фиксированного размера SDF, оптимизация выходит на плато за несколько сотен итераций. Для увеличения предела качества реконструкции нужно повышать разрешение сетки (для удобства мы считаем, что этот шаг отделяет эпохи оптимизации). От того, сколько эпох мы можем себе позволить, сохранив достаточно высокую скорость итерации и приемлемый размер представления, напрямую зависит качество итоговой модели. Для SBS, повышение разрешения заключается в добавлении новых дистанций и получении из одного старого брика восьми новых; каждый содержит столько же дистанций, что и исходный, но представляет в восемь раз меньшее пространство. Был добавлен шаг разреживания сетки: брики добавляются и удаляются таким образом, чтобы итоговое представление состояло из бриков с поверхностью, а также слоя с пустым пространством вокруг них толщиной в один брик.

Примеры реконструкции предложенным методом представлены на \hyperref[fig:exampl]{Рисунке \ref*{fig:exampl}} для моделей Dragon и скана бюста Нефертити. Показаны результаты первой, третьей и седьмой (последней) эпохи, а также исходное изображение.

\begin{figure}[ht]
    \centering
    \includegraphics[width=0.75\textwidth]{collage_sparsedr.pdf}
    \caption{Примеры реконструкции для разных разрешений SBS}
    \label{fig:exampl}
\end{figure}


\endinput